\documentclass[12pt,letterpaper]{article}

\usepackage[utf8]{inputenc}
\usepackage[T1]{fontenc}
\usepackage[english]{babel}
\usepackage{float, enumitem, array, listings}
\usepackage{geometry, amsmath, amssymb, amsthm}
\usepackage{graphicx, xcolor, fancyhdr, fancyvrb}
\usepackage[normalem]{ulem}
\usepackage[hidelinks]{hyperref}

\geometry{margin=2cm, top=2cm, bottom=2cm}

\newcommand{\minic}{
  {M\footnotesize INI\normalsize C}
}

% ========================================
% CONSTANTES DEL ENCABEZADO
% ========================================
\newcommand{\encabezadoIzq}{Compiladores 2027-1}
\newcommand{\encabezadoCen}{Práctica 4}
\newcommand{\encabezadoDer}{\textit{Árboles Sintaxis Abstracta}}

% ========================================
% CONFIGURACION DEL ENCABEZADO
% ========================================
\pagestyle{fancy}
\fancyhf{}

\fancyhead[L]{\encabezadoIzq}
\fancyhead[C]{\encabezadoCen}
\fancyhead[R]{\encabezadoDer}

\fancyfoot[C]{\thepage}

\renewcommand{\headrulewidth}{0.1pt}
\renewcommand{\footrulewidth}{0pt}

% ========================================
% CONFIGURACION Y COLORES LSTLISTINGS
% ========================================
\renewcommand{\lstlistingname}{Código}

\definecolor{links}{HTML}{0D1894}

\definecolor{backcolour}{rgb}{0.93, 0.93, 0.93}
\definecolor{codegreen}{rgb}{0, 0.6, 0}
\definecolor{codegray}{rgb}{0.5, 0.5, 0.5}
\definecolor{codepurple}{rgb}{0.58, 0, 0.82}
\definecolor{keywordcolor}{rgb}{0.5, 0.0, 0.5}
\definecolor{codeorg}{HTML}{DB9400}
\definecolor{linkblue}{HTML}{1A167A}
\definecolor{redcolor}{HTML}{7A1805}
\definecolor{greencolor}{HTML}{0E8F10}
\definecolor{greeblucolor}{HTML}{0E6657}
\definecolor{orgcolor}{HTML}{B86D00}
\definecolor{highlightyellow}{HTML}{FFF176}
\definecolor{boxgray}{HTML}{F2F2F2}
\definecolor{borderblue}{HTML}{1A167A}

% ========================================
% CONFIGURACION LENGUAJES
% ========================================
\lstdefinestyle{modernc}{
    language=C,
    backgroundcolor=\color{backcolour},   
    commentstyle=\color{codegreen},
    keywordstyle=\color{keywordcolor}\bfseries,
    numberstyle=\tiny\color{codegray},
    stringstyle=\color{codepurple},
    basicstyle=\ttfamily\footnotesize,
    breakatwhitespace=false,
    breaklines=true,
    captionpos=b,
    keepspaces=true,
    numbers=left,
    numbersep=5pt,
    showspaces=false,                
    showstringspaces=false,
    showtabs=false,                  
    tabsize=4,
    frame=single,
    rulecolor=\color{black!20},
    frameround=tttt,
    emph={int,uint32_t,char,double,float,unsigned,void,bool,struct,size_t,
          Room, Client, pthread_mutex_t, NULL},
    emphstyle=\color{codeorg}\bfseries
}

\begin{document}

\thispagestyle{empty}
\begin{titlepage}
  \thispagestyle{empty}
  %--------------- COMANDO LINEA----------------->
  \newcommand{\linea}{\rule{\linewidth}{0.5mm}}                 
  \center
  %<<<<<<<<<<<<<<<<<<<<<<<<<<<<<<<<<<<<<<<<<<<

  %Add logos
  \includegraphics[width=0.24\textwidth]{../../unam_logo.png} \hspace{1.5cm}
  \includegraphics[width=0.25\textwidth]{../../fc_logo.png} \\[1cm]
  \textbf{\Large UNIVERSIDAD NACIONAL AUTÓNOMA}\\[3mm]
  \textbf{\Large DE MÉXICO} \\[6mm]
  \textit{\Large Facultad de Ciencias}

  \vfill

  %---------------TÍTULO----------------->
  \linea
  \vfill \vspace{1cm}
  \textbf{\Large Compiladores}\\[1.1cm]
  \textbf{\Large Árboles de sintaxis abstracta}\\[0.9cm]
  \linea \\
  \vfill
  \vspace{0.7cm}
  %Title of the Research
  \textbf{\Large  Práctica 4}\\[6mm]
  %<<<<<<<<<<<<<<<<<<<<<<<<<<<<<<<<<<<<<<<<<<<

  %Team
  \vfill
  \textit{\small Presenta:}\\
  \textbf{\large  \textbf{Ramírez Juárez María Fernanda} - 321204747}\\[3mm]
  \textbf{\large  \textbf{Rojo Peña Manuel Ianluck} - 118005762}\\[3mm]

  %Profesor y Grupo
  \textit{\small Profesor:}\\
  \textbf{Pablo Martínez Yessica Janeth}\\[3mm]
  \textit{\small Ayudante de laboratorio:}\\
  \textbf{Salazar González Pedro Yamil}\\[3mm]
  \small Grupo: 7010, 2027-1\\[3mm]
  \vfill

  %Fecha or Date
  \textit{\small Fecha de entrega:}\\
         {\large 2 de octubre, 2026}
         
\end{titlepage}


%%%%%%%%%%%%%%%%%%%%%%%%%%%%%%%%%%%%%%%%%%
\newpage
%%%%%%%%%%%%%%%%%%%%%%%%%%%%%%%%%%%%%%%%%%

% ================================================================
\section*{Objetivo}
% ================================================================

\begin{quote}
  En esta práctica nuestro objetivo fue implementar la generación de árboles de sintaxis abstracta para el parser
  de \minic ...
\end{quote}


%%%%%%%%%%%%%%%%%%%%%%%%%%%%%%%%%%%%%%%%%%
\bigskip
%%%%%%%%%%%%%%%%%%%%%%%%%%%%%%%%%%%%%%%%%%


% ================================================================
\section*{Evolución respecto de la Práctica 3}
% ================================================================

\begin{lstlisting}[style=modernc, caption={\texttt{fn} implementada}]
\end{lstlisting}

\bigskip

%%%%%%%%%%%%%%%%%%%%%%%%%%%%%%%%%%%%%%%%%%
\bigskip
%%%%%%%%%%%%%%%%%%%%%%%%%%%%%%%%%%%%%%%%%%


% ================================================================
\section*{Diseño del ASA}
% ================================================================


%%%%%%%%%%%%%%%%%%%%%%%%%%%%%%%%%%%%%%%%%%
\bigskip
%%%%%%%%%%%%%%%%%%%%%%%%%%%%%%%%%%%%%%%%%%


% ================================================================
\section*{Construcción desde el Parser}
% ================================================================


\begin{verbatim}

\end{verbatim}


%%%%%%%%%%%%%%%%%%%%%%%%%%%%%%%%%%%%%%%%%%
\bigskip
%%%%%%%%%%%%%%%%%%%%%%%%%%%%%%%%%%%%%%%%%%


% ================================================================
\section*{Propiedad y memoria}
% ================================================================


%%%%%%%%%%%%%%%%%%%%%%%%%%%%%%%%%%%%%%%%%%
\bigskip
%%%%%%%%%%%%%%%%%%%%%%%%%%%%%%%%%%%%%%%%%%


% ================================================================
\section*{Impresión canónica}
% ================================================================


%%%%%%%%%%%%%%%%%%%%%%%%%%%%%%%%%%%%%%%%%%
\bigskip
%%%%%%%%%%%%%%%%%%%%%%%%%%%%%%%%%%%%%%%%%%


% ================================================================
\section*{Diagnósticos y recuperación}
% ================================================================


%%%%%%%%%%%%%%%%%%%%%%%%%%%%%%%%%%%%%%%%%%
\bigskip
%%%%%%%%%%%%%%%%%%%%%%%%%%%%%%%%%%%%%%%%%%


% ================================================================
\section*{Resultados y pruebas}
% ================================================================


%%%%%%%%%%%%%%%%%%%%%%%%%%%%%%%%%%%%%%%%%%
\bigskip
%%%%%%%%%%%%%%%%%%%%%%%%%%%%%%%%%%%%%%%%%%


% ================================================================
\section*{Problemas encontrados}
% ================================================================


%%%%%%%%%%%%%%%%%%%%%%%%%%%%%%%%%%%%%%%%%%
\bigskip
%%%%%%%%%%%%%%%%%%%%%%%%%%%%%%%%%%%%%%%%%%


% ================================================================
\section*{Conclusiones}
% ================================================================


%%%%%%%%%%%%%%%%%%%%%%%%%%%%%%%%%%%%%%%%%%
\bigskip
%%%%%%%%%%%%%%%%%%%%%%%%%%%%%%%%%%%%%%%%%%


% ================================================================
\section*{División del trabajo}
% ================================================================


%%%%%%%%%%%%%%%%%%%%%%%%%%%%%%%%%%%%%%%%%%
\bigskip
%%%%%%%%%%%%%%%%%%%%%%%%%%%%%%%%%%%%%%%%%%


% ================================================================
\section*{Uso de herramientas de Inteligencia Artificial}
% ================================================================


%%%%%%%%%%%%%%%%%%%%%%%%%%%%%%%%%%%%%%%%%%
\bigskip
%%%%%%%%%%%%%%%%%%%%%%%%%%%%%%%%%%%%%%%%%%


% ================================================================

\begin{thebibliography}{99}
\bibitem{ref1}
  ppreference.com. (n.d.). C language. Disponible en \textcolor{links}{\uline{\url{https://en.cppreference.com/c/language}}}

\bibitem{ref2}
  ppreference.com. (n.d.). C11. Disponible en \textcolor{links}{\uline{\url{https://en.cppreference.com/c/11}}}
\end{thebibliography}

\end{document}
